\documentclass[10pt,a4paper]{amsart}
\usepackage[margin=19mm]{geometry}
\usepackage{amsmath,amssymb,mathtools}
\usepackage[scr=rsfs,cal=euler]{mathalfa}
\usepackage{xfrac}
\usepackage[unicode,hidelinks,bookmarks=false]{hyperref}
\newcommand{\ie}{i.e.\ }
\newcommand{\cf}{cf.\ }
\newcommand{\resp}{respectively\ }
\newcommand{\Subsection}{\subsection}
\providecommand{\nobreakdash}{\nobreak\hbox{-}}
\setlength{\emergencystretch}{2em}
\multlinegap0pt
\newcommand{\PGL}{\operatorname{PGL}}
\usepackage{amssymb}
\usepackage{xcolor}
\newtheorem{lemma}{Lemma}
\newtheorem{proposition}{Proposition}
\newcommand{\bl}{\lambda}
\newcommand{\bs}{\sigma}
\newcommand{\ind}{\operatorname{ind}}
\newcommand{\mcM}{\mathcal M}
\newcommand{\sms}{\smallskip}
\title{An analogue of irreducible cuspidal representations for the group $\PGL(2)$ over a two-dimensional local field}
\author{Alexander Braverman, David Kazhdan}
\date{Source: arXiv:2604.11735v2 (CC BY 4.0)}
\begin{document}
\maketitle

\noindent\emph{Source and licence.} An analogue of irreducible cuspidal representations for the group $\PGL(2)$ over a two-dimensional local field. Version arXiv:2604.11735v2 (CC BY 4.0), distributed by arXiv under the Creative Commons Attribution 4.0 International licence, http://creativecommons.org/licenses/by/4.0/, as recorded in the arXiv licence metadata for this exact version and shown on the abstract page https://arxiv.org/abs/2604.11735v2. Full licence notice: \texttt{licenses/arxiv-2604.11735-CC-BY-4.0.txt}.\par\smallskip

\noindent\textit{Editorial scope.} Counted unit: the article's proposition Upsilonm-M (source0.tex lines 2425--2430). The article's complete original proof of it is delivered with it (source0.tex lines 2431--2496, one \texttt{proof} environment of 66 lines). No mathematics is written, repaired or re-derived: the statement and the proof are the article's own lines, in the article's own notation, and no retained line is substituted or re-flowed.  Retained uncounted as the source-local context the proof needs: lines 988--993.  Nothing was omitted from any retained span, so the package carries no omission record.  No retained line contains a citation, so the wrapper carries no bibliography and no bibliography entry.  No retained line contains a figure environment, an image-inclusion command or drawing code, so no image is delivered and the package declares no asset dependency.  Editorial formatting only, in addition: an AMS \texttt{amsart} preamble that restates the article's own declarations and macros used by the retained lines, namely the environments \texttt{lemma}, \texttt{proposition} on separate counters, together with the standard packages those lines need.  The retained environments print in this document as Lemma 1, Proposition 1 in order of appearance, whereas the article prints the same items with the numbering of its own full document.  The complete native source of the article as distributed in the arXiv e-print for this exact version is bundled unchanged as \texttt{sources/198358/source0.tex}.
\begin{lemma}\label{sec}
There is an isomorphism
$$
\bs(m)\simeq\ind_{\widetilde R_m}^{P(O)}(\bl_m).
$$
\end{lemma}

\begin{proposition}\label{Upsilonm-M}
There is a canonical isomorphism of $P_0$-representations
$$
(\tau_m,\Upsilon_m)\simeq(\rho_m,\mcM),\qquad m\geq0.
$$
\end{proposition}

\begin{proof}
We first treat $m=0$.

\sms

Induction in stages gives
$$
\Upsilon_0\simeq\ind_{P'(O)}^{P_0}\widehat\psi.
$$
For $n\geq0$, put $P_0^{\leq n}=O^*\ltimes t^{-n}O$ and let
$$
\Upsilon_0^n=\ind_{P'(O)}^{P_0^{\leq n}}\widehat\psi.
$$
Then
$$
\Upsilon_0\simeq\varprojlim_{n\geq0}\Upsilon_0^n.
$$
As an $F$-variety, the quotient $P_0^{\leq n}/P'(O)$ admits coordinates
$$
F^*\times(t^{-n}O/O).
$$
The character $\widehat\psi$ supplies the usual twisted translation action on the second factor. The residue pairing
$$
(t^{-n}O/O)\times(O/t^nO)\longrightarrow F,\qquad
(z,y)\longmapsto\operatorname{Res}_{t=0}(zy\,dt)
$$
is perfect. Fourier transformation in the second variable therefore gives an isomorphism
$$
\mathcal F_n:\Upsilon_0^n\stackrel{\sim}{\longrightarrow}
\operatorname{Meas}_c\bigl(F^*\times(O/t^nO)\bigr).
$$
The map
$$
F^*\times(O/t^nO)\longrightarrow O^*/O^*(n+1),\qquad
(a,y)\longmapsto a(1+ty),
$$
is an isomorphism of $F$-varieties. Under this isomorphism, the action of $\widehat a\in O^*$ is pushforward under $x\mapsto ax$, while the action of $\widetilde b$, $b\in t^{-n}O$, is multiplication by
$$
x\longmapsto\psi(bx/t).
$$
Thus $\mathcal F_n$ identifies $\Upsilon_0^n$ with $\mcM_n$ as a representation of $P_0^{\leq n}$.
Choose the self-dual additive Haar measures in the finite Fourier transforms. Then restriction from $t^{-n-1}O/O$ to $t^{-n}O/O$ is transformed into pushforward along
$$
O^*/O^*(n+2)\longrightarrow O^*/O^*(n+1).
$$
Consequently the isomorphisms $\mathcal F_n$ are compatible with the transition maps. Passing to projective limits gives
$$
\Upsilon_0\simeq\varprojlim_n\mcM_n=\mcM
$$
as $P_0$-representations.

For general $m\geq0$, Lemma \ref{sec} and induction in stages give
$$
\Upsilon_m\simeq\ind_{\widetilde R_m}^{P_0}(\bl_m),
\qquad \widetilde R_m=O^*(1)\ltimes t^mO.
$$
Conjugation by $\widehat t^{-m}$ identifies $\widetilde R_m$ with $P'(O)$. If $c\in O$ and $b=t^mc\in t^mO$, then
$$
\bl_m(1,b)=\psi(t^{-m-1}b)=\psi(c/t)=\widehat\psi(1,c).
$$
Thus this conjugation identifies the inducing pair $(\widetilde R_m,\bl_m)$ with $(P'(O),\widehat\psi)$. Consequently it identifies the resulting $P_0$-action with the twist
$$
p\longmapsto\rho_0(\widehat t^{-m}p\widehat t^m)=\rho_m(p).
$$
Hence $\Upsilon_m\simeq\mcM$ with the action $\rho_m$, as asserted.
\end{proof}

\end{document}
